%% ===========================================================================
\FloatBarrier
\section{Boundary conditions}
\label{sec:bc}
%% ===========================================================================

The fourth-order character of \eqref{eq:gov1} means that a boundary condition
on the displacement is not enough: the gradient theory requires a second,
higher-order condition on each edge, and the electrical condition decides
which of the two coupling factors the model returns. The mechanical and electrical conditions and coordinate conventions are
specified below.

\FloatBarrier
\subsection{Mechanical and higher-order conditions}
\label{sec:bc_mech}

On a straight edge segment of the rectangular domain, with constant
outward unit normal $\bm{n}$ and tangent $\bm{t}$, stationarity
of \eqref{eq:energy} gives the classical traction
\begin{equation}
  t_i = \left(\sigma_{ij}-\varsigma_{ijm,m}\right)n_j
      - \partial_t\!\left(\varsigma_{ijm}\,t_j n_m\right),
  \label{eq:traction}
\end{equation}
and the higher-order double traction
\begin{equation}
  m_i = \varsigma_{ijm}\,n_j n_m ,
  \label{eq:doubletraction}
\end{equation}
which is conjugate to the normal derivative of the displacement.
The continuum traction $t_i$ has SI units $\si{\pascal}$ and the double
traction $m_i$ has units $\si{\newton\per\meter}$; integrating either
corresponding virtual-work term along a two-dimensional edge yields energy
per unit out-of-plane width, not the total three-dimensional work. The
flexoelectric part of $\varsigma$ in \eqref{eq:tau} is carried by these two
expressions without further change: it contributes $\mu(\nabla\cdot\bm E)n_i$
to \eqref{eq:traction} from the divergence term,
$+\mu\partial_t[(\bm E\cdot\bm n)t_i]$ through the tangential derivative, and
$-\mu(\bm E\cdot\bm n)n_i$ to \eqref{eq:doubletraction}. For piecewise straight edges, tangential integration by parts also
produces corner contributions where the effective traction jumps; these
must be treated consistently with prescribed nodal displacement and slope
conditions. Higher-order surface traction, double-traction and edge-jump
terms appear explicitly in the independent variational treatments of
Mao et al.~\cite{MaoPurohitAravas2016} and Serrao and
Kozinov~\cite{SerraoKozinov2024}. Those references support the requirement
for higher-order boundary data, but they do not validate the particular
two-dimensional edge treatment or replace a three-dimensional benchmark
for the present model. The finite-element weak form includes the flexoelectric term
$\int\mu\nabla N_\phi^{\mathsf T}\nabla(\varepsilon_{xx}+\varepsilon_{yy})$, consistent with the constitutive coupling stated above. The mechanical boundary classes considered are
\begin{equation}
  \text{clamped: } u_i=0,\ u_{i,j}n_j=0;\qquad
  \text{displacement-supported: } u_i=0,\ m_i=0;\qquad
  \text{free: } t_i=0,\ m_i=0 .
  \label{eq:bcmech}
\end{equation}
The middle condition fixes both displacement components but leaves the normal
slope natural; it should not be confused with a conventional beam simple
support. The reported strip simulations prescribe all eight nodal Hermite
quantities on the selected clamped edge, which enforces both $u_i=0$ and its
normal slope there (and their tangential derivatives). Other mechanical
edges are natural/free, except where an external traction or point load is
applied. The displacement-supported class is not a separate simulated case.

A prescribed clamped slope produces a strain-gradient boundary layer: the gradient energy penalises the curvature that the
prescribed slope forces at the edge, so the solution relaxes over a length set
by the window. The window-size dependence is examined against $a/L$;
$a/L\le0.2$ is an assumed input range, not a verified error bound and
not a restriction on the transverse ratio $a_2/h$.

\FloatBarrier
\subsection{Electrical conditions: open and short circuit}
\label{sec:bc_elec}

Two electrical conditions are used. In the \emph{short-circuit} configuration
the electrode pair is grounded, the internal potential degrees of freedom are
condensed, and the mechanical problem has stiffness $K_{sc}$. In the
\emph{open-circuit} configuration one electrode is grounded and the other
floats at a single potential $V$ carrying zero net charge, which is again a
condensation and gives the rank-one update
\begin{equation}
  K_{oc} = K_{sc} + \frac{1}{\beta}\, d\,d^{\mathsf T},
  \label{eq:rankone}
\end{equation}
with $d$ the condensed coupling vector and $\beta$ the condensed dielectric
coefficient. The two coupling measures used in the analysis are defined below. The \emph{mode-based} factor of the $j$-th mode is
\begin{equation}
  k_j^2=\frac{\omega_{oc,j}^2-\omega_{sc,j}^2}{\omega_{oc,j}^2},
  \label{eq:k2mode}
\end{equation}
where $\omega_{oc,j}$ and $\omega_{sc,j}$ are the open- and short-circuit
angular frequencies. This factor depends on the returned eigenvector, while the
\emph{mode-free quasi-static} factor of the electrode pair is
\begin{equation}
  x = \frac{d^{\mathsf T}K_{sc}^{-1}d}{\beta},\qquad
  k^2 = \frac{x}{1+x},
  \label{eq:kstatic}
\end{equation}
which does not require an eigenvector. Equation~\eqref{eq:kstatic} is the
measure used wherever a quantity has to be compared across geometries or
across thicknesses, because it cannot select a different mode in a different
configuration. The following notation distinguishes the coupling measures: $k^2_{\rm qs}=x/(1+x)$ is the mode-free factor of
\eqref{eq:kstatic}; $k^2_{\rm qs,u}=x$ is the same quantity without the coupled
denominator, the one that can be set against $e^{\mathsf T}C^{-1}e/\kappa$;
$k^2_{j}$ is the mode-based factor of the $j$-th mode; and $k^2_{j}(K_{uu})$ or
$k^2_{j}(K_{sc})$ names the matrix whose short-circuit eigenvalues enter it. An
unsubscripted $k^2$ in the results is $k^2_{\rm qs}$. The extensional coupling
factor of the \SI{200}{\micro\meter} by \SI{10}{\micro\meter} strip obtained
from \eqref{eq:kstatic} is $k^2=0.5727$, within $0.5\,\%$ of the constitutive
prediction $0.5752$ obtained independently from the constitutive tensors.

\FloatBarrier
\subsection{Poling axes and Voigt indices}
\label{sec:axes}

The material is tetragonal $\mathrm{BaTiO_3}$ at room temperature, with Voigt
indices $1=x$, $2=y$, $3=z$; the length runs along $x$
(\SI{200}{\micro\meter}), the thickness along $y$ (\SI{10}{\micro\meter}), and
$z$ is out of plane, the neutral axis of the bending and a free surface. The
crystal is anisotropic, so the two poling configurations
do not share a modulus: poling along $x$ pairs the strain $\varepsilon_{xx}$
with $C^{p}_{xx}=\SI{9.5500e10}{\pascal}$, while poling along $y$ pairs the
same strain with $C^{p}_{xx}=\SI{1.6946e11}{\pascal}$, the two being exchanged
by the rotation that takes one poling axis into the other, and the beam moduli
$E_b=1/(\bm C_p^{-1})_{xx}$ are \SI{7.6327e10}{\pascal} and
\SI{1.3544e11}{\pascal}. The electric field runs along the poling axis in
both configurations. The unreduced polar-axis permittivity is
$\kappa_{33}=\SI{4.9583e-10}{\farad\per\meter}$, but the coupled
plane-stress elimination yields the effective value
$\kappa^{p}_{\mathrm{pol}}=\SI{4.9804e-10}{\farad\per\meter}$
across the electrode gap: $\kappa^{p}_{yy}$ with poling along $y$ and
$\kappa^{p}_{xx}$ with poling along $x$. A modulus and a permittivity are
always taken from the same reduced configuration.

The electrode configuration also determines the field direction and gap. The
extensional configuration is poled along $x$ with electrodes on the short
faces and the gap $L=\SI{200}{\micro\meter}$; the bending configuration is
poled along $y$ with electrodes on the long faces and the gap
$h=\SI{10}{\micro\meter}$. The two gaps differ by a factor of twenty.
For uniform open-circuit polarization and equal permittivity, $E=-P/\kappa$
is independent of gap and the voltage scales with gap. At fixed applied
voltage, the field instead scales inversely with gap. The actual coupling
also depends on the strain field, reduced tensors and electrode arrangement. These kinematic and electrical conditions are imposed in the finite-element
formulation.

\FloatBarrier
